\documentclass{professional_resume}

\name{RISHI KIRAN KARNATAKAM}
\address{karnatakamishil@gmail.com $\diamond$ 8328838715 $\diamond$ \href{https://github.com/RishiKarnatakan}{github.com/RishiKarnatakan} $\diamond$ \href{https://linkedin.com/in/rishi-kiran-karnatakan}{linkedin.com/in/rishi-kiran-karnatakan}}

\begin{document}

\rSection{Summary}

\rSection{Experience}

\rSection{Education}
\begin{rSubsection}{NNRESGI}{Dec 2021 - Present}{Bachelor of Technology}{Computer Science and Engineering}
\item GPA: 8.4
\end{rSubsection}
\begin{rSubsection}{KMIT}{Sep 2019 - Jun 2021}{Intermediate, 10+2 equivalent}{}
\item GPA: 9.4
\end{rSubsection}
\begin{rSubsection}{SAGHS}{May 2019}{Secondary School Certificate}{}
\item GPA: 9.8
\end{rSubsection}

\rSection{Skills}
\begin{rSkills}
\item Programming Languages: C, Python
\item Frameworks \& Libraries: Flask, Flutter
\item Databases: MySQL, MongoDB
\item Cloud \& Deployment: Amazon Web Services (AWS), Supabase, Netlify, Google Cloud Platform
\item Machine Learning: TensorFlow, Google Gemini API
\item Tools \& DevOps: Git, GitHub
\end{rSkills}

\rSection{Projects}
\begin{rSubsection}{AI-Powered Resume Builder}{Present}{Supabase, Gemini API, Python, Netlify, Google Cloud Platform}{}
\item Developed an AI-powered resume builder, optimizing content based on job descriptions and user input.
\item Integrated Google Gemini API to intelligently generate and modify LaTeX resume code in real-time.
\item Designed an interactive AI-driven editing environment, similar to Cursor, where users can refine resume content and formatting directly through Gemini's suggestions.
\item Constructed a full-stack application using Supabase for robust user authentication and efficient data storage.
\item Deployed frontend on Netlify and managed backend services on Google Cloud Platform (GCP).
\end{rSubsection}
\begin{rSubsection}{AI-Integrated Plant Disease Detection Mobile App}{Apr 2025}{Flutter, Dart, TensorFlow Lite, MongoDB Atlas}{}
\item Built a mobile app for real-time plant disease detection using camera or gallery images.
\item Integrated a TensorFlow Lite model (Inception V3) for accurate disease prediction with offline support.
\item Enabled local data storage and automatic syncing to MongoDB Atlas when connected online.
\item Designed a clean, modern UI with a sidebar to view history, profile, and AI predictions.
\end{rSubsection}
\begin{rSubsection}{AI-Powered Chess Engine}{Aug 2024}{Python, Tkinter, Chess Library}{}
\item Developed an AI-driven chess engine utilizing minimax and alpha-beta pruning for optimized decision-making.
\item Enhanced strategic depth and move accuracy by 20\%.
\item Achieved a 30\% improvement in move generation speed through alpha-beta pruning.
\end{rSubsection}
\begin{rSubsection}{Cotton Plant Disease Detection and Classification Using Transfer Learning}{Feb 2024}{Python, TensorFlow}{}
\item Developed a machine learning model to classify various diseases on cotton leaves through image processing techniques.
\item Improved disease detection accuracy by 25\% using transfer learning with pre-trained CNN models.
\item Achieved 90\%+ accuracy on test datasets, optimizing performance for agricultural disease diagnosis.
\end{rSubsection}

\rSection{Awards}
\begin{rAward}{National Hackathon on Generative AI}{}{Sep 2024}
\item Achieved 2nd place in a National Level Hackathon focused on Generative AI.
\end{rAward}
\begin{rAward}{Internal Hackathon on Generative AI}{}{Aug 2024}
\item Achieved 3rd place in a college-hosted Hackathon focused on Generative AI.
\end{rAward}
\begin{rAward}{National Hackathon on AR/VR Development}{}{Oct 2023}
\item Ranked in the top 10 at a national-level hackathon focused on AR/VR.
\end{rAward}

\rSection{Certifications}
\begin{rCertification}{AWS Academy Graduate - AWS Academy Cloud Architecting}{AWS Academy}{Sep 2024}
\end{rCertification}
\begin{rCertification}{Python For Data Science}{IIT Madras (NPTEL)}{Jan 2025}
\end{rCertification}
\begin{rCertification}{AICTE Virtual Internship}{Nalla Narasimha Reddy Education Society's Group of Institutions}{2024}
\end{rCertification}
\begin{rCertification}{The Joy of Computing Using Python}{IIT Madras (NPTEL)}{Jul 2023}
\end{rCertification}
\begin{rCertification}{Supervised Machine Learning: Regression and Classification}{Stanford University (Coursera)}{Jun 2023}
\end{rCertification}

\end{document} 